\documentclass[12pt,a4paper]{article}
\usepackage[utf8]{inputenc}
\usepackage[french]{babel}
\usepackage[T1]{fontenc}
\usepackage{geometry}
\geometry{margin=2.5cm}
\usepackage{array}
\usepackage{titlesec}
\usepackage{tabularx}
\usepackage{enumitem}
\usepackage{booktabs}
\usepackage{fancyhdr}

% --- Configuration des entêtes ---
\pagestyle{fancy}
\fancyhf{}
\lhead{Business Plan : Obscurus Anti-Cheat Solutions}
\rhead{Mars 2026}
\cfoot{\thepage}

\title{\textbf{Business Plan : Obscurus Anti-Cheat Solutions}}
\author{Maël - Fondateur \& CTO}
\date{Mars 2026}

\begin{document}

\maketitle

\tableofcontents
\newpage

\section{Genèse et Vision du Projet}
Le projet Obscurus répond à l'obsolescence des solutions anti-triche actuelles face aux menaces émergentes comme le DMA (Direct Memory Access) et l'usage de l'IA par les tricheurs. Obscurus propose une architecture hybride adaptative garantissant l'intégrité des jeux vidéo sans compromettre la stabilité ou la confidentialité des utilisateurs.

\section{Organisation des Ressources Humaines}
L'entreprise est structurée en SARL pour allier expertise technique profonde et gestion rigoureuse :
\begin{itemize}
    \item \textbf{Maël (PDG / CTO) :} Conception du cœur système (C++/ASM), développement du pilote Kernel et architecture des micro-services Cloud (Go/Rust).
    \item \textbf{Anna (Responsable Cybersécurité / RSO) :} Direction des missions offensives, reverse-engineering et veille stratégique sur les forums spécialisés.
    \item \textbf{Joan (Responsable Commercial) :} Développement du portefeuille client B2B, marketing de contenu et relations éditeurs.
    \item \textbf{Camille (Responsable RH \& Admin) :} Gestion financière, conformité légale (RGPD/EULA) et support client.
\end{itemize}

\section{Catalogue des Services Obscurus}

\subsection{Obscurus Cloud \& IA}
Protection non-intrusive basée sur l'analyse de données serveurs.
\begin{itemize}
    \item Analyse comportementale en temps réel via les flux de télémétrie.
    \item Détection par Intelligence Artificielle du "Smoothing" et des comportements non-humains.
    \item Infrastructure à faible latence (< 50ms) déployée sur MySolarCloud.
    \item Tarifs : 1200 € HT / mois pour les titres jusqu'à 1 million de joueurs actifs, avec paliers progressifs au-delà.
\end{itemize}

\subsection{Obscurus Kernel Shield}
Sécurité de haut niveau agissant directement sur le système d'exploitation du joueur.
\begin{itemize}
    \item Pilote Ring 0 interdisant l'accès à la mémoire du jeu par des logiciels tiers.
    \item Protection avancée contre le debugging et le reverse-engineering du moteur de jeu.
    \item Certification WHQL assurant une compatibilité parfaite avec Windows.
    \item 2800 € HT / mois par titre, avec un engagement annuel et des options de personnalisation avancées.
\end{itemize}

\subsection{Audit et Conseil Red Team}
Expertise offensive pour identifier les vulnérabilités avant les tricheurs.
\begin{itemize}
    \item Tests d'intrusion (Pentest) sur le binaire client et les API serveurs.
    \item Identification des failles de logique métier (duplication d'objets, bypass de boutique).
    \item Rapport de remédiation technique avec recommandations de correction de code.
    \item Tarifs : 1200 € HT par mission d'audit, avec des forfaits adaptés pour les studios indépendants et les grands éditeurs.
\end{itemize}

\subsection{Support Premium \& Monitoring 24/7}
Assistance critique pour les lancements et les jeux à fort trafic.
\begin{itemize}
    \item Ligne d'urgence dédiée accessible en permanence (24h/24, 7j/7).
    \item Engagement de réponse (SLA) en moins d'une heure pour les incidents majeurs.
    \item Surveillance proactive de l'infrastructure et ajustement immédiat des filtres de détection.
    \item Tarif : Forfait fixe de 200 € HT / mois.
\end{itemize}

\subsection{Dashboard de Gestion Personnalisable}
Interface de pilotage pour les développeurs et administrateurs.
\begin{itemize}
    \item Visualisation des vagues de bannissements et des statistiques de triche.
    \item Configuration modulaire des niveaux de protection selon les besoins du titre.
\end{itemize}

\section{Plan d'Action Stratégique}
\begin{enumerate}
    \item \textbf{Mois 1-4 :} Développement du MVP Cloud et réalisation des premiers audits Red Team pour financer la R\&D.
    \item \textbf{Mois 5-8 :} Phase de test avec studios partenaires et entraînement des modèles d'IA. Lancement commercial par Joan.
    \item \textbf{Mois 9-15 :} Finalisation du Kernel Shield, certification Microsoft et déploiement des contrats de support premium.
\end{enumerate}

\section{Analyse Financière Prévisionnelle}
Le modèle économique repose sur trois piliers :
\begin{itemize}
    \item Abonnements SaaS indexés sur le volume de joueurs actifs.
    \item Honoraires forfaitaires pour les missions de conseil Red Team.
    \item Revenus récurrents liés aux options de support critique 24/7.
\end{itemize}

\end{document}